% Strict shared typography and layout. The main file must use:
% \documentclass[letterpaper,reqno,11pt]{article}
\usepackage[margin=1.0in]{geometry}
\parskip 1ex
\PassOptionsToPackage{no-math}{fontspec}
\usepackage{fontspec}
\usepackage{xeCJK}
\setCJKmainfont{FandolSong-Regular.otf}[
  BoldFont=FandolSong-Bold.otf,
  ItalicFont=FandolSong-Regular.otf,
  BoldItalicFont=FandolSong-Bold.otf
]
\setCJKsansfont{FandolHei-Regular.otf}[
  BoldFont=FandolHei-Bold.otf,
  ItalicFont=FandolHei-Regular.otf,
  BoldItalicFont=FandolHei-Bold.otf
]
\setCJKmonofont{FandolSong-Regular.otf}[BoldFont=FandolSong-Bold.otf,ItalicFont=FandolSong-Regular.otf,BoldItalicFont=FandolSong-Bold.otf]
\setCJKfamilyfont{kaishu}{FandolKai-Regular.otf}[]
\newcommand{\kaishu}{\CJKfamily{kaishu}}
\xeCJKsetup{AutoFakeBold=false,AutoFakeSlant=false}
\setmainfont{cmunrm.otf}[
  BoldFont=cmunbx.otf,
  ItalicFont=cmunti.otf,
  BoldItalicFont=cmunbi.otf
]
\setsansfont{lmsans10-regular.otf}[
  BoldFont=lmsans10-bold.otf,
  ItalicFont=lmsans10-oblique.otf,
  BoldItalicFont=lmsans10-boldoblique.otf
]
\setmonofont{lmmono10-regular.otf}[
  ItalicFont=lmmono10-italic.otf
]
\usepackage{amsmath,amssymb,amsthm,mathtools,bm}
\usepackage{upgreek,slashed,url,graphicx,amsfonts,xr}
\usepackage[mathscr]{eucal}
\urlstyle{tt}
\numberwithin{equation}{section}
\usepackage{titlesec}
\titleformat*{\section}
  {\normalfont\sffamily\Large\bfseries}
\titleformat*{\subsection}
  {\normalfont\sffamily\large\bfseries}
\titleformat*{\subsubsection}
  {\normalfont\sffamily\normalsize\bfseries}
\pagestyle{plain}
\renewcommand{\contentsname}{目录}
\renewcommand{\refname}{参考文献}
\renewcommand{\indexname}{概念与术语索引}
\renewcommand{\abstractname}{摘要}
\renewcommand{\figurename}{图}
\renewcommand{\tablename}{表}
\renewcommand{\proofname}{证明}
\setcounter{tocdepth}{3}
\setcounter{secnumdepth}{3}
\usepackage[nottoc,notlof,notlot,notindex]{tocbibind}
\usepackage[noautomatic]{imakeidx}
\usepackage[unicode,hidelinks,backref=page]{hyperref}
\usepackage{bookmark}
\usepackage[nameinlink,noabbrev]{cleveref}
\hypersetup{bookmarksnumbered=true,pdfborder={0 0 0},
  pdftitle={规范理论与几何朗兰兹续谈},pdfauthor={Edward Witten}}
\renewcommand*{\backref}[1]{}
\renewcommand*{\backrefalt}[4]{%
  \ifcase #1\relax\space（未被引用）%
  \or\space（引用页：#2）%
  \else\space（引用页：#2）\fi}
\renewcommand*{\backrefsep}{、}
\renewcommand*{\backreftwosep}{、}
\renewcommand*{\backreflastsep}{、}
\makeindex[title=概念与术语索引,intoc,columns=2,options=-s index.ist]
\indexsetup{level=\section*,toclevel=section,firstpagestyle=plain,othercode=\normalsize}
\newcommand{\papercontents}{%
  \clearpage\pdfbookmark[1]{目录}{paper-toc}\tableofcontents\clearpage}
\crefname{equation}{式}{式}
\Crefname{equation}{式}{式}
\crefname{figure}{图}{图}
\Crefname{figure}{图}{图}
\crefname{table}{表}{表}
\Crefname{table}{表}{表}
\crefname{section}{节}{节}
\Crefname{section}{节}{节}
\crefname{subsection}{小节}{小节}
\Crefname{subsection}{小节}{小节}
\crefname{subsubsection}{小节}{小节}
\Crefname{subsubsection}{小节}{小节}
\crefname{theorem}{定理}{定理}
\Crefname{theorem}{定理}{定理}
\crefname{definition}{定义}{定义}
\Crefname{definition}{定义}{定义}
\crefname{lemma}{引理}{引理}
\Crefname{lemma}{引理}{引理}
\crefname{proposition}{命题}{命题}
\Crefname{proposition}{命题}{命题}
\crefname{corollary}{推论}{推论}
\Crefname{corollary}{推论}{推论}
\crefname{remark}{注记}{注记}
\Crefname{remark}{注记}{注记}
\crefname{example}{例}{例}
\Crefname{example}{例}{例}
% Original mathematical notation and Euler/Computer Modern font definitions.
\def\Bbb{\mathbb}
\def\h{\widehat}
\def\bb{b}
\def\cF{{\sf F}}
\def\frak{\mathfrak}
\def\underarrow#1{\vbox{\ialign{##\crcr$\hfil\displaystyle
 {#1}\hfil$\crcr\noalign{\kern1pt\nointerlineskip}$\longrightarrow$\crcr}}}
\def\tr{{\rm tr}} \def\Tr{{\rm Tr}}
\def\gx#1{g_{\hbox{\fvfont #1}}}
\def\N{{\cal N}}
\def\Pic{{\mathrm{Pic}}}
\def\P{p}
\def\pp{{\sf p}}
\def\qq{{\sf q}}
  
\def\hat{\widehat}
% use of underarrow
%A~~~\underarrow{a}~~~B
%

\font\teneurm=eurm10 \font\seveneurm=eurm7 \font\eighteurm=eurm8 \font\fiveeurm=eurm5
\newfam\eurmfam
\textfont\eurmfam=\teneurm \scriptfont\eurmfam=\seveneurm
\scriptscriptfont\eurmfam=\fiveeurm
\def\eurm#1{{\fam\eurmfam\relax#1}}
\font\teneusm=eusm10 \font\seveneusm=eusm7 \font\fiveeusm=eusm5
\newfam\eusmfam
\textfont\eusmfam=\teneusm \scriptfont\eusmfam=\seveneusm
\scriptscriptfont\eusmfam=\fiveeusm
\def\eusm#1{{\fam\eusmfam\relax#1}}
\font\tencmmib=cmmib10 \skewchar\tencmmib='177
\font\sevencmmib=cmmib7 \skewchar\sevencmmib='177
\font\fivecmmib=cmmib5 \skewchar\fivecmmib='177
\newfam\cmmibfam
\textfont\cmmibfam=\tencmmib \scriptfont\cmmibfam=\sevencmmib
\scriptscriptfont\cmmibfam=\fivecmmib
\def\cmmib#1{{\fam\cmmibfam\relax#1}}
\def\FF{{\mathscr F}}
\overfullrule=0pt
\def\G{{\Gamma}}
\def\eG{{\eurm G}}
\def\neg{\negthinspace}
\def\MM{\mathscr{M}}
\def\M{{\MM}}
\def\CP{{\Bbb{CP}}}
\def\L{{\mathcal L}}
\def\U{{\mathcal U}}
\def\CN{{\cal V}}
\def\Y{{\mathcal Y}}
\def\V{{\mathcal V}}
\def\ade{{\mathrm{ad}(E)}}
\def\H{{\mathcal H}}
\def\W{{\mathcal W}}
\def\MH{{{\MM}_H}}
%macros
%
%\def\ppi{{\large\boldsymbol{\mathscr{\pi}}}}
\def\ppi{{\boldsymbol{\mathscr{\pi}}}}
\def\tilde{\widetilde}
\def\bar{\overline}
\def\Z{{\Bbb Z}}
\def\AA{{\eurm A}}
\def\LLL{{\mathfrak L}}
\def\B{{\mathcal B}}
\def\T{{\bf T}}
\def\t{\widetilde }
\def\CO{{\mathcal O}}
\def\O{{\mathcal O}}
\def\F{{\mathcal F}}
\def\FF{{\eusm F}}
\def\A{{\mathcal A}}
\def\CW{{\mathcal W}}
\def\S{{\bf S}}
\def\R{{\Bbb R}}
\def\Lam#1{\Lambda_{#1}}
\def\SIgma{\Sigma}
\def\pf{{\rm Pf ~}}
\font\zfont = cmss10 %scaled \magstep1
\font\litfont = cmr6 \font\fvfont=cmr5
\def\bigone{\hbox{1\kern -.23em {\rm l}}}
\def\ZZ{\hbox{\zfont Z\kern-.4emZ}}
\def\half{{\litfont {1 \over 2}}}
\def\mx#1{m_{\hbox{\fvfont #1}}}
\def\gx#1{g_{\hbox{\fvfont #1}}}
\def\gG{{\cal G}}
\def\lamlam#1{\langle S_{#1}\rangle}
\def\CM{{\cal M}}
\def\C{{\Bbb C}}
\def\G{{\mathcal G}}
\def\Re{{\rm Re ~}}
\def\Im{{\rm Im ~}}
\def\lfm#1{\medskip\noindent\item{#1}}
%$\Bbb{ARZ}$
\def\g{{\frak g}}
\def\CA{{\cal A}}
\def\CF{{\cal F}}
\def\ad{{\rm ad}}
\def\EE{{\mathcal E}}
%\def\PP{\boldsymbol{\mathcal{P}}}
\def\PP{\large{\mathscr P}}
\def\DD{\large{\mathscr D}}
\def\kk{{\mathscr{K}}}
\def\RR{\mathcal{R}}
%\def\tt{{\mathscr{t}}}
\def\ttt{{\bf{t}}}
\def\cc{{\mathscr{C}}}
\def\CZ{?? Z}
\def\D{{\mathcal D}}
\def\rr{{\sf r}}

% ===== 校订译注版新增（2026-11-19）=====
% 译者注以脚注形式随文标出，前缀「译注｜」，与作者的「原注」区分。
\newcommand{\transnote}[1]{\footnote{\textbf{译注｜}#1}}
% 附录：译注汇总与术语对照表
\usepackage{booktabs}
\usepackage{longtable}
\usepackage{enumitem}
\usepackage{xurl}
